\documentclass[11pt]{article}

\usepackage[margin=1in]{geometry}
\usepackage[T1]{fontenc}
\usepackage[utf8]{inputenc}
\usepackage{lmodern}
\usepackage{enumitem}
\usepackage[hidelinks]{hyperref}
\usepackage{titlesec}

\titleformat{\section}{\large\bfseries}{\thesection.}{0.5em}{}
\titleformat{\subsection}{\normalsize\bfseries}{\thesubsection.}{0.5em}{}

\setlist[itemize]{leftmargin=1.25em, itemsep=0.25em, topsep=0.25em}
\setlist[enumerate]{leftmargin=1.5em, itemsep=0.25em, topsep=0.25em}

\title{INJECT and SPECTRA: Updated 2-Month Work Plan}
\author{Sneha Nair and Pauline Barmby\\Western University\\in conjunction with CANDIAPL}
\date{}

\begin{document}

\maketitle

\noindent This plan is tailored to the current star-cluster injection pipeline, with emphasis on rigorous validation, clearer documentation, and a few high-value improvements that should make the project easier to trust, use, and extend.

\section{Goals}

\begin{enumerate}
    \item Increase confidence in scientific correctness.
    \item Make the pipeline easier to run, understand, and reproduce.
    \item Improve notebook and document quality so results are easier to share.
    \item Add targeted enhancements that reduce runtime friction without expanding scope too much.
\end{enumerate}

\section{Priority Areas}

\subsection{Rigorous Testing}

\begin{itemize}
    \item Expand unit tests for light profiles, PSF convolution, catalog generation, and I/O.
    \item Add integration tests for full injection and detection flows.
    \item Add regression tests for known edge cases and failure modes.
    \item Add reproducibility checks for seeded catalog generation and batch runs.
    \item Add multiband consistency tests if multiband workflows are part of the intended use case.
    \item Add performance smoke tests for PSF caching and batch execution.
\end{itemize}

\subsection{Documentation}

\begin{itemize}
    \item Clean up the main README so setup, usage, and expected outputs are easy to find.
    \item Add a short quick start path for notebook users.
    \item Document the batch pipeline, checkpoints, cache settings, and multiband behavior.
    \item Add a testing guide that explains how to run unit, integration, and visual verification tests.
    \item Improve inline notebook narration so the demo notebooks read like a workflow, not just a sequence of cells.
\end{itemize}

\subsection{Useful Improvements}

\begin{itemize}
    \item Add CI-friendly test execution if not already present.
    \item Standardize outputs and file naming for saved catalogs, checkpoints, and plots.
    \item Add lightweight benchmarking scripts or notebooks for PSF caching and end-to-end runtime.
    \item Improve validation around input images, catalog bounds, and configuration values.
    \item Add clearer logging and summary stats for batch runs.
\end{itemize}

\section{Eight-Week Plan}

\subsection{Weeks 1--2: Baseline and Test Foundation}

\begin{itemize}
    \item Audit existing tests and identify the highest-risk gaps.
    \item Add or tighten unit tests around:
    \begin{itemize}
        \item \texttt{mag\_to\_flux} and magnitude/flux scaling.
        \item profile normalization and symmetry.
        \item PSF convolution flux conservation.
        \item catalog generation bounds and seeded reproducibility.
    \end{itemize}
    \item Convert any ad hoc verification into repeatable assertions.
    \item Decide which visual checks should remain manual and which can become automated.
\end{itemize}

\noindent\textbf{Deliverable:} A more reliable test suite that can catch basic scientific and numerical regressions.

\subsection{Weeks 3--4: Integration and Regression Testing}

\begin{itemize}
    \item Add end-to-end tests for single-band injection and detection.
    \item Add tests that cover no-PSF, fallback-PSF, and cache-enabled modes.
    \item Add regression tests for previously observed bugs or fragile behavior.
    \item Add tests for edge cases such as:
    \begin{itemize}
        \item injections near image boundaries,
        \item very faint clusters,
        \item very compact or very extended profiles,
        \item empty or tiny catalogs.
    \end{itemize}
    \item If multiband work matters, add a test that verifies the same catalog is applied consistently across bands.
\end{itemize}

\noindent\textbf{Deliverable:} A test layer that exercises the actual pipeline behavior, not just isolated math.

\subsection{Weeks 5--6: Documentation and Workflow Polish}

\begin{itemize}
    \item Rewrite or expand the main README for:
    \begin{itemize}
        \item installation,
        \item quick start,
        \item notebook usage,
        \item batch execution,
        \item outputs and checkpoints.
    \end{itemize}
    \item Add a concise developer guide for running tests and understanding the repo structure.
    \item Update notebook introductions and final summary cells so each notebook states:
    \begin{itemize}
        \item what it demonstrates,
        \item what input data it expects,
        \item what outputs to inspect,
        \item what limitations remain.
    \end{itemize}
    \item Clarify performance notes for PSF caching and batch execution.
\end{itemize}

\noindent\textbf{Deliverable:} A repo that is easier for a new user to pick up without asking for oral context.

\subsection{Weeks 7--8: Targeted Enhancements and Cleanup}

\begin{itemize}
    \item Add a small benchmarking workflow for before/after comparisons.
    \item Improve logging and run summaries so long batch jobs are easier to monitor.
    \item Tighten config validation and error messages.
    \item Standardize plot and result outputs where that reduces confusion.
    \item Revisit any notebook or script that still feels too manual and reduce repetitive steps.
\end{itemize}

\noindent\textbf{Deliverable:} A more polished and maintainable pipeline with better diagnostics and fewer workflow surprises.

\section{Stretch Ideas}

If time remains after the core work, these are the best next candidates:

\begin{itemize}
    \item CI automation for tests and notebook smoke checks.
    \item A small synthetic benchmark dataset for stable regression testing.
    \item A results summary report that combines completeness curves, detection statistics, and PSF-cache performance.
    \item A lightweight API cleanup pass so the pipeline interface is easier to use from notebooks and scripts.
    \item A short tutorial notebook focused on one canonical end-to-end use case.
\end{itemize}

\section{Success Criteria}

By the end of two months, the project should have:

\begin{itemize}
    \item Stronger test coverage over the pipeline's most important behaviors.
    \item Clearer documentation for new and returning users.
    \item Better reproducibility and fewer manual steps.
    \item At least one measurable improvement in usability, diagnostics, or runtime workflow.
\end{itemize}

\section{Suggested Weekly Rhythm}

\begin{itemize}
    \item Monday: review priorities, test failures, and open questions.
    \item Midweek: implement and validate the highest-value change.
    \item Friday: update docs, run the full relevant test slice, and record what changed.
\end{itemize}

\section{Notes}

\begin{itemize}
    \item Keep validation close to the code that changed.
    \item Prefer repeatable tests over notebook-only verification when possible.
    \item Avoid widening scope unless a test or benchmark shows a real need.
\end{itemize}

\end{document}